\section{奖励函数引导：RL 风格的推导}

\subsection{问题设定：带奖励的采样}

假设我们有一个预训练的扩散模型，其学到的数据分布为 $p_0(x_0)$。同时我们有一个连续可微的奖励函数 $r(x_0)$，当输出 $x_0$ 更符合用户偏好时，$r(x_0)$ 值更高。

\begin{importantbox}{逆问题是奖励函数的特例}
所有逆问题都可以视为奖励函数引导的特例。只需定义：
$$r(x_0) = -\| y - \mathcal{A}(x_0) \|^2$$
即 L2 损失的负值作为奖励函数。当输出 $x_0$ 通过映射 $\mathcal{A}$ 后越接近观测 $y$，奖励越高。
\end{importantbox}

奖励函数引导的适用范围远超逆问题。例如：

\begin{itemize}
    \item \textbf{运动生成}：定义奖励函数使得生成的运动轨迹避免与场景中的物体碰撞
    \item \textbf{分子生成}：定义奖励函数使得生成的分子结构具有特定的对接性质
    \item \textbf{任意可微判别模型}：只要有一个可微的评分模型，就可以将其输出作为奖励函数
\end{itemize}

\subsection{KL 约束下的奖励最大化}

我们希望定义一个新的数据分布 $p_0^*(x_0)$，使得从中采样的 $x_0$ 不仅具有高奖励，还不会偏离原始分布太远。形式化为如下优化问题：

$$p_0^* = \arg\max_{p} \; \mathbb{E}_{x_0 \sim p}[r(x_0)] - \frac{1}{\beta} D_{\mathrm{KL}}(p \| p_0)$$

其中 $\beta > 0$ 控制奖励最大化与分布偏离之间的权衡：
\begin{itemize}
    \item $\beta$ 越大，越侧重于奖励最大化，但可能产生不真实的样本
    \item $\beta$ 越小，越接近原始分布，但奖励可能不够高
\end{itemize}

\begin{knowledgebox}{与 RLHF 的联系}
这一公式形式与 LLM 领域的 RLHF（Reinforcement Learning from Human Feedback）完全一致。在 RLHF 中，$p_0$ 是 SFT 模型的分布，$r$ 是奖励模型的输出，$\beta$ 是 KL 惩罚系数。扩散模型的推理时引导和 LLM 的 RLHF 本质上在解决同一个数学问题。
\end{knowledgebox}

\subsection{最优分布的闭式解}

上述优化问题有闭式解。通过拉格朗日乘子法和变分计算，可以证明最优分布为：

$$p_0^*(x_0) = \frac{1}{Z_0} p_0(x_0) \exp(\beta \cdot r(x_0))$$

其中 $Z_0 = \int p_0(x_0) \exp(\beta \cdot r(x_0)) dx_0$ 是归一化常数（配分函数）。可以看到，新分布在原始分布的基础上，按照 $\exp(\beta \cdot r(x_0))$ 进行重加权——高奖励区域的概率密度被放大，低奖励区域被压缩。

\begin{knowledgebox}{能量模型的视角}
$p_0^*(x_0)$ 的形式恰好是一个能量模型（Energy-Based Model），其能量函数为：
$$E(x_0) = -\log p_0(x_0) - \beta \cdot r(x_0)$$
原始数据分布提供了``先验能量''，奖励函数提供了``条件能量''。$\beta$ 控制了二者的相对权重，类似于统计物理中的逆温度参数。当 $\beta \to 0$ 时，$p_0^* \to p_0$（回到原始分布）；当 $\beta \to \infty$ 时，$p_0^*$ 集中在奖励最大的点附近。
\end{knowledgebox}

\subsection{中间时间步的分布与最优价值函数}

\begin{importantbox}{核心挑战：中间时间步}
扩散模型的生成过程涉及从 $x_T$ 到 $x_0$ 的多步去噪。为了从 $p_0^*$ 采样，我们不仅需要知道 $p_0^*$，还需要知道每个中间时间步的分布 $p_t^*(x_t)$。这要求我们定义并计算最优价值函数 $V^*(x_t, t)$。
\end{importantbox}

类比于最终分布中奖励函数 $r(x_0)$ 的角色，中间时间步的分布可以表示为：

$$p_t^*(x_t) = \frac{1}{Z_t} p_t(x_t) \exp(\beta \cdot V^*(x_t, t))$$

其中 $V^*(x_t, t)$ 是最优价值函数（optimal value function），表示从 $x_t$ 出发、遵循最优策略所能获得的期望未来奖励。

\begin{warningbox}{价值函数没有解析解}
与最终分布不同，中间时间步的最优价值函数 $V^*(x_t, t)$ 没有解析形式。这是因为它涉及对所有可能未来轨迹的期望，无法直接计算。因此必须引入近似。
\end{warningbox}

\subsection{价值函数的 Tweedie 近似}

借鉴 DPS 中的 Tweedie 近似思路，可以将最优价值函数近似为：

$$V^*(x_t, t) \approx r(\hat{x}_0(x_t))$$

即用当前时刻对 $x_0$ 的 Tweedie 估计代入奖励函数。这一近似的含义是：中间时间步的``未来期望奖励''可以用``当前最优估计的即时奖励''来代替。

虽然这是一个相当粗糙的近似，但它具有重要的实用价值：

\begin{itemize}
    \item 不需要训练额外的价值网络
    \item 只需利用已有的扩散模型（提供 Tweedie 估计）和奖励函数
    \item 计算开销仅为一次前向传播（Tweedie 估计）+ 一次奖励函数评估 + 一次反向传播（计算梯度）
\end{itemize}

讲者也提到，理论上可以训练一个专门的神经网络来学习最优价值函数 $V^*(x_t, t)$，但这种方法在实践中往往不够稳定。因此 Tweedie 近似虽然简单，却是目前最实用的方案。

\begin{warningbox}{价值函数近似的层级}
从精确到粗糙，价值函数有以下几种近似方式：
\begin{enumerate}
    \item \textbf{精确计算}：需要对所有未来轨迹求期望，计算不可行
    \item \textbf{训练价值网络}：训练 $V_\phi(x_t, t)$ 近似 $V^*$，需要大量训练数据和时间，且可能不稳定
    \item \textbf{蒙特卡洛估计}：从 $x_t$ 出发采样多条轨迹，取平均奖励。计算量大但更准确
    \item \textbf{Tweedie 单步估计}：$V^*(x_t, t) \approx r(\hat{x}_0(x_t))$，最简单但误差最大
\end{enumerate}
DPS 和大部分推理时引导方法采用第4种，因为它不需要任何额外训练。
\end{warningbox}

\subsection{得到 DPS 等价公式}

在上述近似下，新分布的分数函数变为：

$$\nabla_{x_t} \log p_t^*(x_t) = \underbrace{\nabla_{x_t} \log p_t(x_t)}_{\text{原始分数}} + \underbrace{\beta \cdot \nabla_{x_t} r(\hat{x}_0(x_t))}_{\text{期望未来奖励的梯度}}$$

\begin{importantbox}{DPS 的 RL 解释}
将奖励函数取为 $r(x_0) = -\| y - \mathcal{A}(x_0) \|^2$，上式立即退化为 DPS 的公式。因此，DPS 可以从两个角度理解：
\begin{itemize}
    \item \textbf{贝叶斯视角}：条件分数 = 边际分数 + 条件似然梯度
    \item \textbf{RL 视角}：最优策略分数 = 原始策略分数 + 期望未来奖励梯度
\end{itemize}
两者在数学上完全等价，但 RL 视角提供了更一般的框架，奖励函数不限于逆问题。
\end{importantbox}

\subsection{本章小结}

通过 KL 约束的奖励最大化，我们可以定义一个新的目标分布。利用 Tweedie 近似估计最优价值函数后，可以推导出与 DPS 完全等价的引导公式。RL 视角为推理时引导提供了更一般的理论框架。

